\begin{lexlean}{SkipLogic.TaskGraph}
\useglossary{lexlean.std.bool@1.2.0}
\useglossary{lexlean.std.nat@1.2.0}
\useglossary{skiplogic.titles@1.0.0}
\importmodule{SkipLogic.Parameters}
\title{Skip logic lock free dynamic task graph scheduler}
\begin{semanticmodule}
structure{
  name: TaskGraphScheduleConfig,
  type_parameters: [],
  parameters: [],
  fields: [
    {name: maxTaskNodesCapacity, type: Nat},
    {name: workerThreadsCount, type: Nat},
    {name: maxDependencyDepth, type: Nat},
    {name: isLockFreeScheduler, type: Bool}
  ]
}

definition{
  name: defaultTaskGraphConfig,
  parameters: [],
  result: TaskGraphScheduleConfig,
  body: (record TaskGraphScheduleConfig {maxTaskNodesCapacity := 256, workerThreadsCount := 8, maxDependencyDepth := 16, isLockFreeScheduler := true})
}

definition{
  name: isSchedulerLockFree,
  parameters: [c : TaskGraphScheduleConfig],
  result: Bool,
  body: (project c isLockFreeScheduler)
}

definition{
  name: isDagAcyclicBounded,
  parameters: [c : TaskGraphScheduleConfig],
  result: Bool,
  body: ((project c maxDependencyDepth) <=? 64)
}

definition{
  name: canSatisfyDependency,
  parameters: [depRemaining : Nat],
  result: Bool,
  body: (depRemaining =? 0)
}

theorem{
  name: task_graph_lock_free_sound,
  parameters: [],
  statement: ((isSchedulerLockFree (defaultTaskGraphConfig)) = true),
  proof: reflexivity{}
}

theorem{
  name: task_graph_acyclic_bounded,
  parameters: [],
  statement: ((isDagAcyclicBounded (defaultTaskGraphConfig)) = true),
  proof: reflexivity{}
}

theorem{
  name: dependency_satisfaction_sound,
  parameters: [],
  statement: (((canSatisfyDependency 0) && (not (canSatisfyDependency 1))) = true),
  proof: reflexivity{}
}
\end{semanticmodule}
\end{lexlean}
